\documentclass[a4paper,11pt]{article}
\usepackage{exam-style-842-unit}
\examinehead{842 信息化概论 · 第 05 单元练习册}{练习三\quad 大数据与 Hadoop}
\begin{document}

\unitcover{842 信息化概论 · 单元练习册}{第 05 单元\quad 云计算与大数据}{练习三\quad 大数据与 Hadoop（共 100 分）}{约 75 分钟}{大数据的由来、定义与 3V 特征；大数据总体架构；Hadoop 的组成与作用（HDFS、YARN、MapReduce、HBase、Hive、Spark）；MapReduce 原理与执行流程；大数据分析与可视化、数据挖掘；大数据问题与对策。}

\begin{knowmap}
\textbf{由来：}Internet、智能手机、物联网每时每刻产生海量、非结构化、实时的数据 $\rightarrow$ 2013 年中国数据总量超过 0.8ZB（相当于 8 亿 TB）；存储单位按 1024 进位：bit $\rightarrow$ Byte $\rightarrow$ KB $\rightarrow$ MB $\rightarrow$ GB $\rightarrow$ TB $\rightarrow$ PB $\rightarrow$ EB $\rightarrow$ ZB $\rightarrow$ YB $\rightarrow$ DB $\rightarrow$ NB\\[2pt]
\textbf{定义：}狭义大数据是用现有一般技术难以管理的大量数据的集合，特征为 Volume（数据量）、Variety（种类）、Velocity（速度）3V；广义大数据还包括存储处理分析这些数据的技术，以及能从中获得实用意义和观点的人才与组织\\[2pt]
\textbf{Hadoop：}支持数据密集型分布式应用的 Apache 2.0 开源软件框架，根据 Google 的 MapReduce 与 Google 档案系统（GFS）论文制作；5 个优点为高可靠性、高扩展性、高效性、高容错性、低成本\\[2pt]
\textbf{核心组件：}HDFS（分布式存储，NameNode 管元数据、DataNode 存数据）$\cdot$ YARN（资源管理与调度，教材仅在 Spark 处提及）$\cdot$ MapReduce（分布式计算，映射＋化简）$\cdot$ HBase（建立在 HDFS 之上的数据库）$\cdot$ Hive（数据仓库工具）$\cdot$ Spark（存储器内运算框架）\\[2pt]
\textbf{执行流程：}MapReduce 的输入分片 $\rightarrow$ Map 阶段局部处理 $\rightarrow$ Shuffle 阶段按 key 分组排序 $\rightarrow$ Reduce 阶段汇总输出（Shuffle 教材未展开，按通行资料补充）\\[2pt]
\textbf{生态与问题：}从离线计算到实时计算（Storm）、从数据库到数据挖掘（Mahout、MLlib）；隐私是大数据最严重的问题，还面临数据归属、使用权与利用限度、数据质量、人才等问题
\end{knowmap}

\begin{keyterms}
\textbf{大数据}：所涉及的数据量规模巨大到无法通过人工在合理时间内被截取、管理、处理并整理，从而成为人类所能解读的形式的信息。\\
\textbf{狭义大数据}：用现有的一般技术难以管理的大量数据的集合，具有 Volume、Variety、Velocity 3 个特征。\\
\textbf{广义大数据}：包括具备 3V 特征而难以管理的数据、对这些数据进行存储处理分析的技术，以及能获得实用意义和观点的人才和组织。\\
\textbf{Volume（数据量）}：指用现有技术无法管理的数据量，目前基本指从几十 TB 到几 PB 的数量级。\\
\textbf{Variety（种类）}：数据的种类，既包括结构化数据，也包括网站日志、文本、位置信息、传感器、图片和视频等非结构化数据。\\
\textbf{Velocity（速度）}：数据产生和更新的频率，如 POS 数据、点击流数据、高峰时每秒数千条的推文。\\
\textbf{HDFS}：Hadoop 分布式文件系统，位于 Hadoop 底层，采用主/从结构存储集群中所有存储节点上的文件。\\
\textbf{NameNode}：HDFS 主节点，管理文件系统的元数据；实际数据由 DataNode 存储。\\
\textbf{MapReduce}：Hadoop 的分布式计算框架，合并映射（Mapping）与化简（Reducing）两种经典函数，以可靠容错的方式并行处理大数据集。\\
\textbf{Hive}：建立在 HDFS 之上的数据仓库工具，由 Facebook 贡献。\\
\textbf{HBase}：建立在 HDFS 之上的数据库，提供数据库的相关功能（教材 4.2.3）。\\
\textbf{Spark}：开源计算框架，最初由加州大学伯克利分校 AMPLab 开发，采用存储器内运算，适合机器学习等迭代计算。\\
\textbf{YARN}：Hadoop 的资源管理与作业调度层；教材仅在 Spark 管理方式处提到“Hadoop YARN”，未系统介绍，按通行资料补充。\\
\textbf{数据挖掘}：从大量数据中提取隐含的、事先未知的、潜在有用信息的过程；DM 与数据库中的知识发现（KDD）本质含义相同。
\end{keyterms}

\begin{contrast}{狭义大数据 / 广义大数据（每行依次为：范围·核心内容·意义）}
\citem{狭义大数据}{只针对数据本身}{用现有一般技术难以管理的大量数据集合，具有 3V 特征}{强调数据规模、种类与产生速度带来的管理难题}
\citem{广义大数据}{数据＋技术＋人才与组织}{3V 数据、存储处理分析技术（Hadoop、NoSQL、机器学习等）、数据科学家等人才与组织}{更全面地反映大数据的全方位特征，兼具技术性与社会性}
\end{contrast}

\begin{contrast}{HDFS / MapReduce / Spark（每行依次为：定位·核心机制·特点）}
\citem{HDFS}{分布式存储系统，Hadoop 的底层}{主/从结构：NameNode 管元数据，DataNode 存实际数据，数据多副本}{可靠存储、可处理超大文件、流式访问、可运行于廉价商用机器集群}
\citem{MapReduce}{分布式计算框架}{把任务分解为 Map 局部处理与 Reduce 汇总，中间结果落盘}{批处理框架、可靠容错、可用 Java 编写，降低分布式编程门槛}
\citem{Spark}{开源分布式计算框架}{存储器内运算，数据尽量驻留内存并反复查询}{比 Hadoop MapReduce 内存运算快约 100 倍、硬盘运算快约 10 倍，适合机器学习迭代计算}
\end{contrast}

\begin{contrast}{映射（Map）/ 化简（Reduce）（每行依次为：操作含义·扑克牌例子中的对应·作用）}
\citem{映射}{对集合里的每个目标应用同一个操作}{每个玩家数自己手中黑桃的张数}{把大任务拆成可并行的小任务，机器越多算得越快}
\citem{化简}{遍历集合中的元素来返回一个综合的结果}{把所有玩家报出的数字相加}{把并行结果汇总为最终结论}
\end{contrast}

\qtypename{一、单项选择题}{10 题，每题 3 分，共 30 分}
\q 教材 4.2.2 概括的狭义大数据三个特征不包括（\quad）
\choicesline{A. Volume}{B. Variety}{C. Velocity}{D. Value}

\q 所涉及的数据量规模巨大到无法通过人工在合理时间内被截取、管理、处理并整理的信息，称为（\quad）
\choicesline{A. 大数据}{B. 数据仓库}{C. 数据库}{D. 数据字典}

\q 据统计，2013 年中国产生的数据总量超过（\quad）
\choicesline{A. 0.8ZB}{B. 8.5ZB}{C. 80PB}{D. 8EB}

\q 按 1024 的进率计算，1ZB 等于（\quad）EB。
\choicesline{A. 1024}{B. 1000}{C. 512}{D. 2048}

\q Hadoop 是根据 Google 公司的哪两项技术（论文）制作而成的（\quad）
\choices{A. MapReduce 与 Google 档案系统（GFS）}{B. TCP/IP 与 HTTP}
{C. 关系模型与 SQL}{D. 神经网络与深度学习}

\q HDFS 集群中，负责管理文件系统元数据的是（\quad）
\choicesline{A. NameNode}{B. DataNode}{C. JobTracker}{D. Hive}

\q HDFS 中的任意一份数据，理论上至少存放（\quad）副本。
\choicesline{A. 1 份}{B. 2 份}{C. 3 份以上}{D. 10 份以上}

\q 教材中作为数据仓库工具、由 Facebook 贡献的 Hadoop 子项目是（\quad）
\choicesline{A. Hive}{B. HBase}{C. Mahout}{D. Storm}

\q 与 Hadoop MapReduce 相比，Spark 的核心优势是（\quad）
\choices{A. 使用存储器内运算，计算速度快}{B. 不需要任何集群即可运行}
{C. 只能处理结构化数据}{D. 不支持机器学习算法}

\q 教材认为大数据面临的最严重问题是（\quad）
\choicesline{A. 隐私问题}{B. 存储成本问题}{C. 网络带宽问题}{D. 显示器分辨率问题}

\qtypename{二、多项选择题}{4 题，每题 4 分，共 16 分；每题有两个或两个以上正确项}
\q 教材 4.2.1 所述企业大数据的来源包括（\quad）
\choicesline{A. 物联网}{B. Internet 查询数据}{C. 社会化数据}{D. 单一结构化报表}

\q 广义大数据包括（\quad）
\choices{A. 具备 Volume、Variety、Velocity 特征而难以管理的数据}{B. 对这些数据进行存储、处理、分析的技术}
{C. 能够通过分析这些数据获得实用意义和观点的人才和组织}{D. 仅指关系数据库中的结构化数据}

\q 教材 9.6.13 指出，数据挖掘（DM）与数据库中的知识发现（KDD）涉及的范围包括（\quad）
\choicesline{A. 企业数据}{B. 商业数据}{C. Web 数据}{D. 仅限科学实验数据}

\q 下列属于 Hadoop 子项目的有（\quad）
\choicesline{A. Hive}{B. Pig}{C. Zookeeper}{D. 微信}

\qtypename{三、判断题}{6 题，每题 2 分，共 12 分；正确的画“√”，错误的画“×”并写出改正后的表述}
\q HDFS 采用主/从体系结构，NameNode 管理元数据，DataNode 存储实际数据。\quad（\quad）\quad 改正：\fixblank

\q MapReduce 只能用于实时计算，不能进行批处理。\quad（\quad）\quad 改正：\fixblank

\q Spark 采用存储器内运算，运算速度比 MapReduce 快。\quad（\quad）\quad 改正：\fixblank

\q 大数据时代个人数据都是匿名的，不存在隐私问题。\quad（\quad）\quad 改正：\fixblank

\q Hadoop 是以 Apache 2.0 许可协议发布的开源软件框架。\quad（\quad）\quad 改正：\fixblank

\q Hive 是分布式锁设施，Zookeeper 是数据仓库工具。\quad（\quad）\quad 改正：\fixblank

\qtypename{四、配对题}{6 对，每对 1 分，共 6 分；把 A 栏与 B 栏对应内容配对}
\q 将下列 Hadoop 组件与其主要作用配对，把 B 栏字母填入括号。
\par\vspace{0.2em}
\noindent\begin{tabular}{@{}p{0.44\textwidth}@{\hspace{1.2em}}p{0.44\textwidth}@{}}
（1）HDFS \hfill（\quad） & A. 数据仓库工具，由 Facebook 贡献 \\
（2）MapReduce \hfill（\quad） & B. 分布式文件系统，由 NameNode 与 DataNode 组成主/从结构 \\
（3）Hive \hfill（\quad） & C. 分布式锁设施，提供类似 Google Chubby 的功能 \\
（4）Zookeeper \hfill（\quad） & D. 分布式计算框架，合并映射与化简两种经典函数 \\
（5）Spark \hfill（\quad） & E. 采用存储器内运算的开源计算框架 \\
（6）Mahout \hfill（\quad） & F. 提供可扩展机器学习经典算法实现的子项目 \\
\end{tabular}\par\vspace{0.3em}

\qtypename{五、简答题}{3 题，每题 8 分，共 24 分}
\q 什么是大数据？写出教材的狭义定义与广义定义，并说明狭义大数据的三个特征。\anslines{2.4cm}

\q 简述 MapReduce 的原理与作用，并说明它为什么能提高计算速度。\anslines{2.4cm}

\q 简述 Hadoop 平台的核心组成：HDFS、YARN、MapReduce、HBase、Hive、Spark 各自解决什么问题、彼此是什么关系。\anslines{2.4cm}

\newpage
\qtypename{六、综合题}{1 题，共 12 分}
\q 某农业集团汇集四类数据：（1）田间传感器与农机作业数据；（2）气象与遥感影像；（3）电商平台用户评论与交易记录；（4）园区监控视频。集团计划建设农业大数据平台。请回答：
（1）按教材对大数据来源的划分，指出上述数据分别属于哪类来源，并用 Volume、Variety、Velocity 说明其“大”在何处。（4 分）
（2）为该平台设计“采集—存储—处理—分析—应用”的总体架构，并说明 HDFS、MapReduce、Spark、Hive 可分别承担哪个环节。（4 分）
（3）分析该平台在数据隐私、数据质量与人才方面面临的问题及对策，并说明分析结果可用什么方式呈现给管理者。（4 分）\anslines{5cm}

\end{document}
